\section{Benchmarking procedure}\label{sec:procedure}

\subsection{Compute and software}
The pool of 2026 releases and the map--model correlation of the candidates (Section~\ref{sec:datasets}) were computed in free Google Colab notebooks; every run of a method and all scoring ran on free Kaggle notebooks. No paid service was used. A GPU session has two NVIDIA \GpuName{} GPUs (\GpuMemory{} each), \CpuCores{} CPU cores and \RamGb{}~GB of RAM; a CPU session has the same cores and memory and no GPU. Phenix \PhenixVersion{} was installed with the author's academic licence (the installer is not redistributed); US-align and every method were installed from pinned commits (Table~\ref{tab:pins}). Each session wrote the output of \texttt{nvidia-smi}, the number of cores, the memory, the package versions and the hashes of its input files into its result zip. The notebooks are generated by scripts of the repository and were tested against stand-ins of the methods before they were run; the sessions that produced results are listed in Table~\ref{tab:sessions}. A result zip is unpacked by a script that keeps text files only (coordinate files stay in the notebook outputs).

\subsection{Order of work}
(1)~The entries were selected and frozen, and the AF3 predictions were made (Section~\ref{sec:datasets}). (2)~A pilot of every method on the smallest entry (26DK) fixed versions and commands before any production run. (3)~The four non-MICA methods ran on all 20 entries in three batch sessions. (4)~MICA ran in three stages: Merizo domain splitting and docking in CPU sessions, then the network and the model building in GPU sessions. (5)~All models were scored in separate CPU sessions with one notebook code, using the protocol of Section~\ref{sec:metrics}, which was written down on 24~October, before any MICA model existed. (6)~After the first version, CryoAtom ran on all \NEntries{} entries in a GPU session of its own, and CryoZeta ran on as many entries as the free GPU time allowed in four sessions; their models were scored in CPU sessions with the same code (Section~\ref{sec:added}). No entry was dropped or replaced after its first run, and no parameter was changed for a single map.

\subsection{Rules for the runs}
The rules below were written in the project log before the sessions that used them started; Table~\ref{tab:failures} lists every failure with its reason.
\begin{itemize}
\item \textbf{One run per method and entry.} A run that ends without the model file counts as failed, whatever the reason (crash, out-of-memory kill, timeout). Failed runs stay in every table and are scored as zero in a sensitivity analysis (Section~\ref{sec:results}).
\item \textbf{Run limit.} 3~h for every run of the baselines; for MICA, 30~min for program~1, a limit for each docking call (Section~\ref{sec:dockprocedure}) and 2~h (GPU) or 3~h (CPU) for program~3. Reaching a limit is a failure. The spec allowed 11~h per run; it was shortened because a weekly GPU quota of about 30~h could not carry hung runs. CryoAtom and CryoZeta, added later, had 10~h and 6~h per run (Section~\ref{sec:added}).
\item \textbf{Two programs at a time.} A Kaggle GPU session has two GPUs; the baselines, CryoAtom and CryoZeta included, ran one program on each. Their run times are therefore wall-clock times with another program active, and only indicative. E3-CryoFold and EModelX(+AF) never ran at the same time, because both allocate arrays that grow with the square of the number of candidate atoms and ran out of RAM together in the pilot. A baseline that the system killed (exit code $-9$) while the other GPU was busy was run once more, alone; a kill while alone, a crash or a timeout is a failure at once.
\item \textbf{Cut by the end of a session.} A run that the end of a session interrupted is not a result; it was run again in a later session.
\item \textbf{GPU memory.} A CUDA out-of-memory error of CryoAtom or CryoZeta on the 15~GB GPU is a failed run of the method on this hardware and not a problem of the installation (Section~\ref{sec:added}); it stays in the tables and does not stop the notebook's guard, which pauses a method for the rest of a session after two other crashes in a row within an hour.
\item \textbf{Order.} The baselines took the entries in the seeded order with the two strata interleaved, so that any prefix is balanced; nothing was dropped or left out because of a result.
\end{itemize}

\subsection{MICA in three stages}\label{sec:dockprocedure}
\paragraph{Docking (program~2).}
The AF3 prediction of every chain was split into domains by Merizo (program~1; \DockDomainsThree{} domains in the 20 entries). Docking was run three times, always on CPU sessions of at most 12~h, and the three passes are all reported (Table~\ref{tab:docking}).
\emph{Pass~1} called Phenix \texttt{dock\_in\_map} with \texttt{dock\_with\_mr=False quick=True} and at most 45~min per domain, followed (if the quick search had not run out of time) by a thorough search of at most 60~min; it docked \DockDockedOne{} of \DockDomainsOne{} domains, with a median local correlation of \DockCcMedOne{}. \emph{Pass~2} was meant to search at coarser resolution, but Phenix~\PhenixVersion{} rejected the parameter that had been added (a True/False parameter where a number was passed), so the first two attempts stopped after one second and only the call of pass~1 ran again: \DockDockedTwo{} domains. This was an error in the set-up of the pass, found in the logs and recorded in \texttt{LOG.md}; it is the reason for a third pass. \emph{Pass~3}, the set-up that counts, runs the whole docking, the search and the rigid-body refinement, at 6~\AA{} (parameter \texttt{resolution=6.0}, at most 40~min per domain) and at 10~\AA{} if that fails or runs out of time (25~min). It docked \DockDockedThree{} of \DockDomainsThree{} domains, with a median local correlation of \DockCcMedThree{}. That value is not comparable with the \DockCcMedOne{} of pass~1, because it is measured at 6~\AA{} after a coarse refinement; I did not check the placements against the deposited models, and the only evidence that they help is indirect (Section~\ref{sec:dockresults}). By entry, \DockModeFull{} entries have every domain docked (\emph{full}), \DockModePartial{} have some (\emph{partial}) and \DockModeNone{} has none (\emph{none}; 9VI4).

\paragraph{Program~3 on a GPU.}
Program~3 of MICA deletes and rewrites its own output folders, so the notebook runs it for one entry at a time. A guard switches program~3 off for the rest of a session after two failures in a row, and the entries that wait then are marked \emph{deferred} and run in a later session. The first GPU session used MICA's own batching and failed with a GPU out-of-memory error on the maps with many windows (Section~\ref{sec:mica}); from then on \texttt{batch\_threshold} was raised to $10^6$ at run time, which was decided and logged before the rerun, and the notebook stops if the patch does not apply exactly once. In the first session \NRerunMica{} entries (\RerunMicaList{}) failed; they ran again with the patch. This is an exception to the rule of one run per entry, which the spec allows only for runs killed by the platform. The first failure stays in Table~\ref{tab:modes}, and the main analysis uses the rerun; if the reruns counted as failures, MICA would have \number\numexpr\NModelMica-\NRerunMica\relax{} models (Section~\ref{sec:results}). Three entries have no MICA model after all sessions: 26DK and 9VI4 were killed by the system when the model-building step of MICA, which runs on one CPU core, used up the 31~GB of RAM (about 33\,000 and 32\,000 candidate atoms, as counted in the program logs), and 9OC0 reached the 2-hour limit of program~3 while it was filling the last gap of its last chain. Nothing else was running when 26DK and 9VI4 were killed.

\paragraph{A second MICA run without docked input.}
To get an indication of what the docked input is worth, a second GPU session ran program~3 from the same saved state with the docked AF3 models removed, so that the AF3 channels of the network are zero. It covers the \DockEffectN{} entries that the first GPU session built (Section~\ref{sec:dockresults}); the equivalent sessions for the later entries were not run, because the GPU quota did not allow it.

\subsection{Scoring and analysis}
Every model file that a method wrote was scored in CPU sessions with the same notebook code. The scored file of each method is fixed in Table~\ref{tab:methods}; further kinds of output (for example MICA's initial C$\alpha$ model) were scored as well and are reported in the Appendix, but are not used in the main comparison. The analysis script reads the status tables of all sessions and the scores, takes the latest result of every method and entry (for MICA, only from the final docking set-up), and writes the tables, the figures and the numbers quoted in the text (as \LaTeX{} macros), so that the numbers of the results are not typed in by hand; the few figures quoted from logs (for example the candidate atoms above) are marked as such. Its quality checks stop on a method that is unknown, an entry that is not in the frozen list, a score for a run that did not finish, and more than one primary model per method and entry; the check on entries that have both an \emph{ok} and a \emph{failed} result flags exactly the \NRerunMica{} MICA reruns above.
